\section{Plone Project Analysis} % {{{
\label{sec:Plone Project Analysis}

\marginpar{\includegraphics[width=\marginparwidth]{plone/logo}}

\subsection{Introduction} % {{{
\label{sub:Introduction}

Plone is a \ac{CMS} built on top of the Zope application server which provides
the core functionality layer \cite{Aspeli2005,PloneFaq,PloneWhatIsPlone}. It is
mostly written in Python and available for all major platforms. Plone is
licenced under the \ac{GPL} and according to Ohloh, Plone is one of the top
\unit[2]{\%} of open source projects worldwide \cite{PloneOhlohFactoids}.
Furthermore it is used by organizations and companies such as NASA, Amnesty
Internation, Nokia and others. It can be used for a wide range of websites such
as internal websites, blogs, groupware and other. The project is known for it's
very good security track record, high usability, extensibility and flexibility.
The project is split into Plone Core and Collective. This analysis however will
only focus on Plone Core as it is the heart of the project.

% }}}

\subsection{History} % {{{
\label{sub:History}

In 1999, Alexander Limi and Alan Runyan started the Plone Project creating a
usability layer on top of the application server and content management
framework Zope \cite{Aspeli2005,PloneFaq}. The first version then was released
two years later in 2001 \cite{PloneReleases}. It was quickly picked up by many
people and a community around the Plone project began to emerge. In 2004 Plone
2.0 was released which made Plone more configurable and enhanced the
possibility of adding modules to the \ac{CMS}. At the same time, the Plone
foundation was founded, which to this day has ownership rights over the Plone
project and trademarks. The next major Plone version was released in 2007 under
the name Plone 3. The nowadays used and developed Plone 4 version was released
in 2010.

% }}}

\subsection{Community} % {{{
\label{sub:Community}

The Plone project consists of a large user and developer community. While there
are about 300 people working on Plone Core, there is a much bigger community
working on Plone Collective projects
\cite{Aspeli2005,PloneOhlohFactoids,PloneCommunityProcesses}.

Most of the communication is done through the Plone mailing lists whereas the
plone-developers mailing list is in place to discuss the development and future
of Plone Core. There are also lots of other mailing lists available for almost
every aspect concerning the project. However, there are also large forums and
\ac{IRC} channels available where one can find support.

The first Plone conference was held in New Orleans in 2003
\cite{PloneConferences}. Since then the yearly Plone conference attracted more
people every year and was held in Europe as well as in the United States.

\begin{figure}[htbp]
  \centering
  \includegraphics[width=0.8\textwidth]{plone/commits_by_author}
  \caption{Monthly activity of the most active Plone Developers}
\end{figure}

The Plone community is also very eager to organize so called Sprints
\cite{PloneSprints}. A Sprint is a focused development sessions where
developers meet in person and try to enhance a certain part of the Plone
project. It normally consists of about three to five days in which one person
leads the sessions and tracks the activities and the development. There are
many Sprints available worldwide every year.

Due to the large group of developers of Plone, the project is split into
several teams, where each team is responsible for a certain part of the Plone
project such as releases, the development of plone but also security,
infrastructure, usability and other
\cite{PloneFounders,PloneReleaseManagers,PloneFrameworkTeam,PloneContribute}.

\begin{figure}[htbp]
  \centering
  \includegraphics[width=0.8\textwidth]{plone/commits_by_year}
  \caption{Yearly overview of commits to Plone}
\end{figure}

\begin{description}

  \item[Core Developer] This role is defined to have access to the Plone Core
    codebase. To become a member of this role one must have established a track
    of contributions to the Plone project. This is not limited to code, but can
    also be usability enhancements or other contributions. It is also highly
    recommended to attend a conference or sprint in order to become
    acquainted with the other developers. Finally one have to request an
    account. Informally, Godefroid Chapelle, Hanno Schlichting and Martin
    Aspeli serve as gatekeepers for this role. Currently there this role has
    375 members.

  \item[Framework Team] This role is responsible for the decision of which code
    gets into a Plone release. They do this by recommending code for inclusion
    to the Plone release manager team. They are however not responsible for
    adding new features or patching bugs, they only drive the \ac{PLIP}
    process. By accepting or rejecting \acp{PLIP} they gather information from
    the community and finally recommend them to the release manager.
    Furthermore all significant changes must go through the Framework Team. The
    team consists of a small group of people and are chosen arbitrarily.
    Currently the team consists of 16 people who are Alec Mitchell, Andreas
    Zeidler, Calvin Hendryx-Parker, David Glick, Elizabeth Leddy, Erik Rose,
    Laurence Rowe, Martijn Pieters, Matthew Wilkes, Raphael Ritz, Rob Gietema,
    Ross Patterson, Tom Lazar, Craig A. Haynal, Danny Bloemendaal, Yiorgis
    Gozadinos.

  \item[Release Managers] Each major Plone release is done by a release manager
    who is responsible for continuing releases in the specific series. The
    release manager gets named by the Plone Foundation board and is often a paid
    position \cite{PlonePaidReleaseManager}. He works closely with the Framework
    team to ensure improvement and continuing development of Plone. The current
    release managers are Alec Mitchell, Eric Steele, Hanno Schlichting, Stefan H.
    Holek, Wichert Akkerman, Alan Hoey.

  \item[Plone Founders] The Plone founders Alexander Limi and Alan Runyan hold
    this role. While it does not give any further rights, both still have a
    powerful voice in the community and help with decisions where the community
    is not able to take one. In some cases they also take a decision by
    themselves, overruling the community's decision.

\begin{figure}[hbtp]
  \centering
  \includegraphics[width=\textwidth]{plone/punchcard}
  \caption{Time based view on commits of contributors}
\end{figure}

\end{description}

% }}}

\subsection{Release Process} % {{{
\label{sub:Release Process}

Since the Plone 4.2 release in 2011 the Plone project changed it's release
process in order to implement a fixed release cycle
\cite{PloneFixedReleaseCycle}. Every six months, there will be a new Plone
release. Additionally the Plone project sets a feature freeze date two months
before the release at which no new \ac{PLIP} will get accepted unless it has
been completely reviewed by that date.

\begin{figure}[htbp]
  \centering
  \includegraphics[width=\textwidth]{plone/releases}
  \caption{Major and feature releases of Plone}
\end{figure}

The Plone project uses a major, minor and micro versioning scheme
\cite{PloneReleaseProcess,PloneCommunityProcesses}. Major releases break
backwards compatibility but may provide an upgrade path. Minor or feature
releases normally provide new features and bugfixes but retain backwards
compatibility. Micro releases are only bug fix releases and do not contain new
features. Additionally, each major and minor release has several preceding
releases. Alpha releases are basically snapshots of the current development
process. Once a beta release gets published, no new feature is allowed to enter
the codebase. Only bug fixes and non-invasive changes are allowed. A release
candidate is published afterwards and only if problems are found, further
release candidates follow. If no show stopper bug comes up, the last release
candidate is considered to be the final release and gets released on the
previously set release date. After that date only bug fix releases are allowed
for this series. The development starts with the next major or minor release.

% }}}

\subsection{Development} % {{{
\label{sub:Development}

Most of the development of the Plone project is handled over the mailing list
and bug trackers of the project \cite{PloneContribute,PloneCommunityProcesses}.
Minor patches and contributions often get committed directly to the repository
while bigger changes and features have to go through the Plone Improvement
Proposal process
\cite{PlonePLIPProcess,PloneCommunityProcesses,PlonePLIPLifecycle}. This
process is based on the Python project's Python Enhancement Proposal process.
New feature and big changes have to be written down in a \ac{PLIP} which
contain a description of the feature or problem, a solution to it and a working
implementation. A \ac{PLIP} is always owned by at least one person who is
responsible for it.

A \ac{PLIP} mostly comes up from a discussion on the plone-developers mailing
list. The community will give feedback on the idea and possible solutions. If
the feedback can be considered as positive, the creator of the idea will be
asked to write a \ac{PLIP}. However one can also start directly with the
creation of a \ac{PLIP} and get the community feedback during the upcoming
phases.

The author of the \ac{PLIP} is always considered as the main contact point for
future improvements or discussions. From here the \ac{PLIP} life cycle begins
\cite{PlonePLIPLifecycle}.

\begin{figure}[htbp]
  \centering
  \includegraphics[width=0.8\textwidth]{plone/commits_by_month}
  \caption{Amount of commits per month of Core Contributors}
\end{figure}

\begin{figure}[htbp]
  \centering
  \begin{tikzpicture}[auto]
   \node[rectangle, rounded corners, draw=black, text width=6.5em, minimum height=2em, text centered]
     (draft) {Draft};
   \node[below=0.5cm of draft, rectangle, rounded corners, text width=6.5em, minimum height=2em, text centered]
     (dummy) {};
   \node[below=0.5cm of dummy, rectangle, rounded corners, draw=black, text width=6.5em, minimum height=2em, text centered]
     (withdrawn) {Withdrawn};

   \node[right=1cm of draft, rectangle, rounded corners, draw=black, text width=6.5em, minimum height=2em, text centered]
     (proposal) {Proposal};
   \node[below=0.5cm of proposal, rectangle, rounded corners, draw=black, text width=6.5em, minimum height=2em, text centered]
     (approved) {Approved};
   \node[below=0.5cm of approved, rectangle, rounded corners, draw=black, text width=6.5em, minimum height=2em, text centered]
     (rejected) {Rejected};

   \node[right=1cm of proposal, rectangle, rounded corners, draw=black, text width=6.5em, minimum height=2em, text centered]
     (review) {Review};
   \node[below=0.5cm of review, rectangle, rounded corners, draw=black, text width=6.5em, minimum height=2em, text centered]
     (final review) {Final review};
   \node[below=0.5cm of final review, rectangle, rounded corners, draw=black, text width=6.5em, minimum height=2em, text centered]
     (accepted) {Accepted};

   \path[draw, -latex] (draft) -- (proposal);
   \path[draw, -latex] (draft) -- (withdrawn);
   \path[draw, -latex] (proposal) -- (approved);
   \path[draw, -latex] (approved.east) -| ($(approved.east) + (0.5cm,0)$) |-(review.west);
   \path[draw, latex-latex] (review) -- (final review);
   \path[draw, -latex] (final review) -- (accepted);

   \path[draw, latex-latex] ($(proposal.south) - (0.5cm,0)$) -- ++(0,-0.3cm) --
   ++(-1.5cm,0) |- ($(rejected.west) + (0,0.1cm)$);
   \path[draw, -latex] ($(rejected.west) - (0,0.2cm)$) -- ($(withdrawn.east) - (0,0.2cm)$);
  \end{tikzpicture}
  \caption{Possible paths of the status of \acp{PLIP}}
\end{figure}

\begin{description}

  \item[Draft] A \ac{PLIP} always starts a draft, where the author summarizes
    his ideas and tries to come up with a solution. It is not ready for a
    proposal but the discussion can already start and new solutions or idea can
    find their way into the \ac{PLIP}.

  \item[Proposal] Once the author thinks that the \ac{PLIP} is finished and
    will get accepted he can submit it at any time. The Framework Team will
    then comment on any newly submitted or updated \ac{PLIP}. It will be
    reviewed and if the Framework Team decides that the \ac{PLIP} won't get
    accepted for Plone Core, the author is encouraged to improve the \ac{PLIP}
    and resubmit it at any later point in time.

\begin{figure}[hbtp]
  \centering
  \includegraphics[width=0.8\textwidth]{plone/authors_by_month}
  \caption{Amount of distinct authors of Plone Core over time}
\end{figure}

  \item[Approved] As soon as the \ac{PLIP} was approved, the author of the
    \ac{PLIP} should start with the development of the \ac{PLIP}. A member of
    the Framework Team will get assigned to the \ac{PLIP} and assists the
    author with any questions and help he might need.

  \item[Rejected] If the \ac{PLIP} was rejected by the Framework Team, in most
    cases the \ac{PLIP} gets implemented as an amendment and later submitted
    again as an improved \ac{PLIP}.

  \item[Review] Once the development is complete, the author of the \ac{PLIP}
    can forward the implementation to the Framework Team. Two members will then
    review the proposal along with two additional developers the author can
    choose. The reviews have to be in written form and available for a all
    members of the community. This phase should not take longer than two weeks
    and all critical bugs have to be fixed before going on. If there are too
    many changes for the available time, the \ac{PLIP} will be pushed to the
    next release.

  \item[Final Review] Once all major concerns have been solved, the original
    reviewers will review it one more time and confirm that all issues have
    been resolved and no new issues have been found. If issues have been found,
    the \ac{PLIP} goes back to the Review state.

  \item[Accepted] If the \ac{PLIP} passes the Final Review phase, it can be
    merged and will be available in the next Plone release.

  \item[Withdrawn] An author can of course always decide to withdraw the
    \ac{PLIP} if he thinks, that the \ac{PLIP} is no longer needed or obsolete.

\end{description}

% }}}

% }}}
